\begin{theorem}[The round's graded-agreement composition refines its specification]
  \label{thm:aba-gbca-composition-by-abdy-refines-specification}
  \lean{PLTS.ABA.GBCA.ByABDY.refinesSpecification}\lean{PLTS.ABA.GBCA.ByABDY.specificationRelation_transition}\lean{PLTS.ABA.GBCA.ByABDY.composition_projects}\lean{PLTS.ABA.GBCA.specificationLabelMap}\lean{PLTS.ABA.GBCA.specificationOverRoundAlphabet}\lean{PLTS.ABA.GBCA.ByABDY.specificationRelation_init}\lean{PLTS.ABA.GBCA.ByABDY.refinesSpecification_failAct}\lean{PLTS.ABA.GBCA.ByABDY.specificationCorruptionAct}\leanok
  \uses{def:aba-gbca-instance-by-abdy, def:aba-gbca-algorithm-by-abdy, def:aba-gbca-spec, def:forwardsim}
  Each round's composition is forward-simulated by that round's GBCA specification instance, read over the round's interface, along a relation that certifies the exclusion set and the grade by receipt patterns, fires an exclusion on demand, holds at the initial states and survives the corruption broadcast: \leandecl{PLTS.ABA.GBCA.ByABDY.refinesSpecification}.
\end{theorem}

\begin{proof}\leanok
  \uses{def:aba-gbca-instance-by-abdy, def:aba-gbca-algorithm-by-abdy, def:aba-gbca-spec, def:forwardsim, def:weakstep}
  See the Lean proof of \leandecl{PLTS.ABA.GBCA.ByABDY.refinesSpecification}, which reads a transition of the composition as one of the algorithm at \leandecl{PLTS.ABA.GBCA.ByABDY.composition_projects} and matches that transition at \leandecl{PLTS.ABA.GBCA.ByABDY.specificationRelation_transition}.
\end{proof}
